% --- 概要 --- 
\begin{abstract}
    ここに論文の概要（アブスト）を記述する．
    通常このセクションでは研究の目的，手法，主な結果，結論を簡潔にまとめる．
\end{abstract}

\pagenumbering{arabic} % ページ番号を 1, 2... にリセット
\setcounter{page}{1}   % 強制的に1ページ目にする



% --- 序論 ---
\begingroup
    \let\clearpage\relax        % 改ページ命令を無効化
    \let\cleardoublepage\relax  % 偶数奇数ページ調整も無効化
    \section{序論}
\endgroup
ここに序論を書く．この研究の背景や重要性，既存の手法について説明する．そしてこの論文で何をやったのかを説明をする．
本文で各内容を明確にするため，適当に流さず「目的（何を解決するためにこの論文を書いた？）」や「動機（なぜこの論文を書きたいと思った？）」，「意義（この論文で誰がハッピーになる？）」などを
数式をなるべく用いずに，わかりやすく簡潔に示す大事な章である．

読者（審査員）はここを読んで「面白そうか」「読む価値があるか」を判断する．